\chapter{Концевой допуск минимальной симплектической достройки}
\label{ch:v8-horizontal-phase-minimal-symplectic-completion-endpoint-admission-gate}

Предыдущий гейт локализовал шестимерный радикал в избытке трёх слабых
дублетов отрицательного гиперзаряда. Теперь различим два утверждения:
существование абстрактного симплектического завершения и его происхождение
из уже фиксированных концевых блоков конечного оператора.

Минимально сбалансированное представление имеет вид
\begin{equation}
 \mathcal T_{26}
 =H_+^{\oplus4}\oplus H_-^{\oplus4}
 \oplus C_+\oplus C_-\oplus S_0^{\oplus4},
 \qquad \dim_{\mathbb C}\mathcal T_{26}=26,
 \qquad \dim_{\mathbb R}\mathcal T_{26}=52.
 \label{eq:v8-minimal-symplectic-completion-decomposition}
\end{equation}
Здесь типы $H_\pm,C_\pm,S_0$ те же, что в предыдущем гейте; новых
калибровочных зарядов не вводится.

В порядке блоков
$(H_+^{\oplus4},H_-^{\oplus4},C_+,C_-,S_0^{\oplus4})$
предъявим форму
\begin{equation}
 \Omega_0=
 \begin{pmatrix}
  0&I_8&0&0&0\\
  -I_8&0&0&0&0\\
  0&0&0&I_3&0\\
  0&0&-I_3&0&0\\
  0&0&0&0&J_4
 \end{pmatrix},
 \qquad
 J_4=\begin{pmatrix}0&1&0&0\\-1&0&0&0\\0&0&0&1\\0&0&-1&0\end{pmatrix}.
 \label{eq:v8-minimal-symplectic-standard-form}
\end{equation}
Для четырнадцати точных образующих калибровочного действия машинно
проверено
\begin{equation}
 \rho(X)^{\mathsf T}\Omega_0+\Omega_0\rho(X)=0,
 \qquad \Omega_0^{\mathsf T}=-\Omega_0.
 \label{eq:v8-minimal-symplectic-gauge-invariance}
\end{equation}
Кроме того,
\begin{equation}
 \Omega_0^2=-I_{26},\qquad
 \Omega_0^{-1}=-\Omega_0,\qquad
 \det\Omega_0=1,\qquad \rank\Omega_0=26.
 \label{eq:v8-minimal-symplectic-nondegeneracy}
\end{equation}
Следовательно, формальный симплектический носитель действительно существует.

Он по-прежнему не каноничен. Общая инвариантная форма задаётся матрицей
$A\in M_4(\mathbb C)$ между слабодвойственными пространствами, одним
цветным коэффициентом и формой $B\in\Lambda^2\mathbb C^4$ на нейтральных
синглетах. Поэтому
\begin{equation}
 \dim\bigl(\Lambda^2\mathcal T_{26}^*\bigr)^G
 =4\cdot4+1+\binom42=23,
 \label{eq:v8-minimal-symplectic-form-family}
\end{equation}
а невырожденность означает $\det A\neq0$, ненулевой цветной коэффициент и
$\operatorname{Pf}B\neq0$. Перестановка двух копий $H_-$ даёт вторую
отличную инвариантную форму полного ранга.

Главная проверка относится к происхождению новых направлений. Текущий
концевой носитель содержит только одну независимую копию $H_+$, тогда как
симплектическая форма требует четыре:
\begin{equation}
 m_+^{\rm endpoint}=1,qquad m_+^{\rm symp}=4,qquad
 \dim\operatorname{coker}\pi_+=3,qquad
 \Delta\dim_{\mathbb C}=2\cdot3=6.
 \label{eq:v8-minimal-symplectic-endpoint-deficit}
\end{equation}
Совпадение калибровочного типа делает эти направления допустимыми как
внешние бозонные котангенциальные переменные, но не увеличивает кратность
образа внутренних флуктуаций. В существующей конечной тройке отсутствуют
три независимых концевых блока, из которых они могли бы возникнуть.

Для одного коммутирующего поля запрет сохраняется,
\begin{equation}
 \Phi^{\mathsf T}\Omega_0\Phi=0,
 \qquad
 \Phi^{\mathsf T}\Omega_0\Psi\not\equiv0
 \quad\text{для независимых }\Phi,\Psi.
 \label{eq:v8-minimal-symplectic-two-field-pairing}
\end{equation}
Точный базисный свидетель имеет значение $e_1^{\mathsf T}\Omega_0e_9=1$.
Тем самым завершение создаёт нужную двухполевую свёртку, но лишь после
введения невыведенного второго набора переменных.

Стандартная симплектическая структура на удвоенном колчане согласуется с
котангенциальной интерпретацией пространства представлений
\cite{v8-minimal-symplectic-crawley-boevey,v8-minimal-symplectic-harada-wilkin}.
Однако классификация конечных спектральных троек требует отдельно предъявить
новые блоки оператора и проверить первое условие и вещественную структуру
\cite{v8-minimal-symplectic-krajewski}.

\begin{theorem}[Формальный допуск минимального завершения]
Представление $\mathcal T_{26}$ допускает калибровочно-инвариантную
невырожденную кососимметрическую форму. Её ранг равен $26$, определитель
равен единице, а пространство всех инвариантных форм имеет размерность
$23$. Два независимых поля допускают ненулевую симплектическую свёртку.
\end{theorem}

\begin{nogocriterion}[Концевой no-go]
Минимальная достройка не является внутренней флуктуацией уже заданного
концевого носителя: ему недостаёт трёх копий $H_+$, то есть шести
комплексных направлений. Их можно добавить как новую котангенциальную
архитектуру, но нельзя считать выведенными из текущего родителя. Новые
фермионы не следуют логически из абстрактного бозонного удвоения; если же
требовать реализации новыми концевыми состояниями, аномалии, первое условие
и лёгкий спектр должны быть проверены заново.
\end{nogocriterion}

\begin{protocol}
\begin{enumerate}
 \item формальная комплексная размерность: \textbf{$26$};
 \item формальная вещественная размерность: \textbf{$52$};
 \item инвариантных кососимметрических форм: \textbf{$23$};
 \item невырожденный представитель: \textbf{предъявлен};
 \item его ранг и определитель: \textbf{$26$ и $1$};
 \item каноническая форма выбрана калибровочной симметрией: \textbf{нет};
 \item ненулевая двухполевая свёртка: \textbf{да};
 \item дефицит концевой кратности $H_+$: \textbf{$3$};
 \item новых комплексных направлений: \textbf{$6$};
 \item происхождение из текущих внутренних флуктуаций: \textbf{нет};
 \item физический подъём горизонтальной фазы: \textbf{не получен};
 \item следующий гейт: \textbf{допуск котангенциального удвоенного
 колчана как нового родителя}.
\end{enumerate}
\end{protocol}

Машинный сертификат сохранён в
\path{s2t_v8_horizontal_phase_minimal_symplectic_completion_endpoint_admission_gate_results.json}.

\begin{thebibliography}{9}
\bibitem{v8-minimal-symplectic-crawley-boevey}
W.~Crawley-Boevey,
\emph{Poisson structures on moduli spaces of representations},
J. Algebra 325 (2011), 205--215; arXiv:math/0502301.
\bibitem{v8-minimal-symplectic-harada-wilkin}
M.~Harada, G.~Wilkin,
\emph{Morse theory of the moment map for representations of quivers},
Geom. Dedicata 150 (2011), 307--353; arXiv:0807.4734.
\bibitem{v8-minimal-symplectic-krajewski}
T.~Krajewski,
\emph{Classification of Finite Spectral Triples},
J. Geom. Phys. 28 (1998), 1--30; arXiv:hep-th/9701081.
\end{thebibliography}